% =====================================================================================================================
% Supplementary material of the IEEE paper (compile on its own: Overleaf Menu > Main document = supplementary.tex, pdfLaTeX).
% Floats are numbered in source order: Tables S1-S8 and Figures S1-S5 are cited by number in main.tex, so keep their order.
% =====================================================================================================================
\documentclass[journal]{IEEEtran}

\usepackage[T1]{fontenc}
\usepackage[utf8]{inputenc}
\usepackage{amsmath,amssymb,amsfonts}
\usepackage{graphicx}
\usepackage{booktabs}
\usepackage{array}
\usepackage{tabularx}
\usepackage{url}
\usepackage{xcolor}
\usepackage{siunitx}
\usepackage{adjustbox}
\usepackage[hidelinks]{hyperref}

\setlength{\emergencystretch}{2em}
\renewcommand{\topfraction}{0.9}\renewcommand{\dbltopfraction}{0.9}\renewcommand{\textfraction}{0.07}
\renewcommand{\floatpagefraction}{0.7}\renewcommand{\dblfloatpagefraction}{0.7}
\setcounter{topnumber}{4}\setcounter{dbltopnumber}{4}\setcounter{totalnumber}{6}
\renewcommand{\thesection}{S\arabic{section}}
\renewcommand{\thesubsection}{S\arabic{section}.\arabic{subsection}}
\renewcommand{\thetable}{S\arabic{table}}
\renewcommand{\thefigure}{S\arabic{figure}}

\newcommand{\todo}[1]{\textcolor{red}{[TODO: #1]}}
\newcommand{\safefig}[2][\linewidth]{%
  \IfFileExists{#2}%
    {\includegraphics[width=#1]{#2}}%
    {\fbox{\parbox[c][2.4cm][c]{0.92\linewidth}{\centering\todo{missing figure #2}}}}}
\newcommand{\restable}[2][\columnwidth]{\begin{adjustbox}{max width=#1}\footnotesize\input{tables/#2.tex}\end{adjustbox}}

\begin{document}
\begin{center}
{\Large\bfseries Supplementary Material}\\[0.5em]
{\large Lesion-Faithful Bone Suppression from Chest CT Projections: Counterfactual Auditing and Gains for Focal Lesion Detection and
Classification in Chest Radiographs}
\end{center}
\medskip

This document gives the analysis notes and the secondary tables and figures of the paper. Tables and figures of the main paper are cited
there as ``Table~I'' and ``Fig.~1'', those of this document as ``Table~S1'' and ``Fig.~S1''.

% =====================================================================================================================
\section{Analysis Notes and Corrections}
\label{sup:notes}
This section records, in one place, every change to the analysis that was made after it started, with its reason. None changed the
verdict of a hypothesis.

\subsection{Study I}
\begin{itemize}
\item \emph{Choice of encoder.} A replacement rule declared before the final runs moved ConvNeXt-V2-T into the place of the planned
 SwinV2-T. At two seeds (development data) RAD-DINO reached $A=0.046$ and bone-region PSNR$_{\mathrm{cal}}$ 21.92\,dB and
 ConvNeXt-V2-T $A=0.050$ and 22.01\,dB; both met the rule, and the tie-break declared in advance (mean absorption within 0.005, then the
 higher bone-region PSNR) chose ConvNeXt-V2-T, the only encoder that also passed the screening rule (Table~\ref{stab:screen}), so excluding
 RAD-DINO would not have changed the choice. With the third seeds the absorption gap is 0.0055, just outside the band; the choice was not
 reopened and both encoders are reported in the main paper.
\item \emph{Encoder learning rate.} ConvNeXt-V2-T diverged at the full learning rate (fp16 overflow) and was rerun with an encoder learning
 rate $0.25\times$ that of the decoder, an extension of the rule written for vision transformers.
\item \emph{Development use of SZCH.} After the v1 ablation the SZCH pairs were looked at, so SZCH became a development set for decisions
 about the final recipe. While the recipe was being chosen, screening tools printed held-out PSNR and SSIM next to the development
 numbers for earlier models and screening runs; no rule, table or decision used them, and the tool was changed to print development
 numbers only.
\item \emph{JSRT overlap map.} JSRT has no bone reference. The protocol named the consensus of the bone maps of ours and of the reference
 model, but the audit took the overlap from the reference model's bone map alone (a secondary analysis).
\item \emph{BDCXR-3257 duplicate grouping.} After the protocol of the external validation was fixed, near-duplicate images were grouped
 by their correlation at 64\,px ($r\ge0.93$, 48 pairs) so that no group spans two folds; this was decided from the images alone before
 any model ran.
\item \emph{Corrections found in an independent check.} After the first evaluation of the locked sets, an independent check of the
 statistics code found a resampling error that had narrowed interaction intervals and detection metrics that depended on the order of tied
 scores by up to 0.00055; both were corrected and all analyses rerun. Outputs from cloud and local GPUs differ by at most 0.0024\,dB.
\item \emph{Execution.} Some evaluations ran on cloud GPUs instead of the local GPU with identical code; the lung masks were exported once
 so that every model is audited at the same sites.
\end{itemize}

\subsection{Study II}
\begin{itemize}
\item \emph{Gradient accumulation.} SwinV2-B did not fit a 16\,GB GPU at batch size 32 and was screened with four micro-batches per step
 (the same effective batch). The other encoders ran unchanged.
\item \emph{Encoder learning rate in the screen.} ConvNeXt-V2-T stayed at chance at the planned learning rate. The three encoders
 without batch normalisation were screened again with the encoder learning rate reduced fourfold, decided before the other screen results
 were seen, on validation data only.
\item \emph{Holm families.} Before any test evaluation we fixed how H6 and H7 would be corrected if two encoders went forward, and which
 models H10 uses. With one encoder the families are those of the main paper.
\item \emph{Repeated NIH runs.} One account's NIH soft-tissue array differed at the byte level from the other five copies. The five runs
 trained on it (soft seeds 0, 3, 4; early seeds 2, 3) were repeated on a consistent copy before the analysis. With the superseded runs,
 the early and soft seed ensembles would have scored 0.8279 and 0.8253 instead of 0.8273 and 0.8243; this check entered no test.
\item \emph{Early fusion.} The contrast between early fusion and the radiograph alone (0.8273 against 0.8241 on NIH) was not one of the
 tested hypotheses and is descriptive.
\end{itemize}

% =====================================================================================================================
\section{Additional Tables}
\label{sup:tables}
Tables~\ref{stab:szch}--\ref{stab:rib} give every method and seed of the main fidelity, audit and rib-benchmark tables,
Table~\ref{stab:seeds} the single seeds of the final models, Table~\ref{stab:qp} the exact-truth audit on the Qure.ai test split,
Table~\ref{stab:amplitude} the amplitude sweep of the audit, Table~\ref{stab:screen} the encoder screen and Table~\ref{stab:v1} the v1
hypothesis H1, Table~\ref{stab:per_class} the per-class AUROC of the fusion study and Table~\ref{stab:cls_uq} the uncertainty of the
classifiers.

\begin{table*}[!t]
\centering
\caption{Held-out SZCH-X-Rays pairs (51): bone-region and whole-image metrics for every method and seed, calibrated (primary) and raw.
HF ratio: high-frequency energy of the non-bone region relative to the dual-energy soft-tissue image (1 = equal).}
\label{stab:szch}
\restable[\textwidth]{tmi_s_szch}
\end{table*}

\begin{table}[!t]
\centering
\caption{Counterfactual lesion audit on VinDr-RibCXR: mean absorption of every method and seed with image-cluster bootstrap intervals
(lower is better).}
\label{stab:audit}
\restable{tmi_s_audit_ribcxr}
\end{table}

\begin{table}[!t]
\centering
\caption{Rib benchmark on VinDr-RibCXR for every method (metrics only).}
\label{stab:rib}
\restable{tmi_s_rib}
\end{table}

\begin{table*}[!t]
\centering
\caption{Single-seed results of the final models. The ensembles average seeds 0 to 2; SwinV2-T seed 3 is the reserve.}
\label{stab:seeds}
\restable[\textwidth]{app_seeds}
\end{table*}

\begin{table*}[!t]
\centering
\caption{Exact-truth audit on the Qure.ai test split (430 projections; lesions added to the soft component): absorption, lesion contrast
ratio of the soft estimate and bone MAE. The SwinV2-T twin row averages the four seeds of the dose-response arm ($K=4$); the other trained
models use seeds 0--2.}
\label{stab:qp}
\restable[\textwidth]{app_qp_audit}
\end{table*}

\begin{table}[!t]
\centering
\caption{Audit amplitude sweep on the 450 SZCH sites (development data).}
\label{stab:amplitude}
\restable{app_amplitude}
\end{table}

\begin{table}[!t]
\centering
\caption{Encoder screen at \num{100000} samples, twin weight 16 (development data).}
\label{stab:screen}
\restable{app_encoder_screen}
\end{table}

\begin{table*}[!t]
\centering
\caption{v1 hypothesis H1 on synthetic LIDC-IDRI test CTs (exact truth; three seeds per arm): absorption of the insertion arms against the
recipe arm and the non-inferiority of the bone MAE.}
\label{stab:v1}
\restable[\textwidth]{v1_h1}
\end{table*}


\begin{table*}[!t]
\centering
\caption{Fusion study, per class (H8): test AUROC of the radiograph and hybrid seed ensembles and the paired DeLong contrasts of the hybrid against the
radiograph and against channel stacking; Holm within data set over classes and both contrasts.}
\label{stab:per_class}
\restable[\textwidth]{v3_per_class}
\end{table*}

\begin{table*}[!t]
\centering
\caption{Classification uncertainty. Top: conformal risk control of the false-negative rate on NIH (nominal 0.10) and class-conditional
(Mondrian) sets on BDCXR-3257 (nominal coverage 0.90). Bottom: selective classification, area under the risk--coverage curve with the
entropy of the seed ensemble alone and with the suppressor's spread added (lower is better).}
\label{stab:cls_uq}
\restable[\textwidth]{v3_conformal}
\par\smallskip\noindent{\footnotesize\textit{Selective classification}}\par\smallskip
\restable[\textwidth]{v3_selective}
\end{table*}

% =====================================================================================================================
\section{Additional Figures}
\label{sup:figures}
Fig.~\ref{sfig:reliability} shows the per-image coverage of the bone-image intervals, Fig.~\ref{sfig:bone_unc} their recalibration, Fig.~\ref{sfig:seeds} every run of the fusion study,
Figs.~\ref{sfig:conformal} and~\ref{sfig:selective} the classification uncertainty, Fig.~\ref{sfig:ig}
Integrated-Gradients maps and Fig.~\ref{sfig:pointing} the pointing game per class.

\begin{figure*}[!t]
\centering
\safefig[0.78\textwidth]{figures/reliability.pdf}
\caption{Left: per-image coverage of the conformal intervals on the Qure.ai test split and on the held-out SZCH-X-Rays pairs (bars:
means; dashed: nominal 0.90). Right: per-image Spearman correlation of ensemble spread and absolute error on SZCH-X-Rays (line: mean).}
\label{sfig:reliability}
\end{figure*}

\begin{figure*}[!t]
\centering
\safefig[0.9\textwidth]{figures/fig_bone_uncertainty.pdf}
\caption{Bone-image intervals on the 51 held-out SZCH-X-Rays pairs, recalibrated on the 24 development pairs (H9, post hoc): per-image
coverage, reliability and sparsification.}
\label{sfig:bone_unc}
\end{figure*}

\begin{figure*}[!t]
\centering
\safefig[0.8\textwidth]{figures/fig_seed_spread_densenet121.pdf}
\caption{Fusion study: every single run next to the seed ensembles of each configuration.}
\label{sfig:seeds}
\end{figure*}

\begin{figure*}[!t]
\centering
\safefig[0.8\textwidth]{figures/fig_conformal_sets_densenet121.pdf}
\caption{Conformal sets per class (conformal risk control on NIH ChestX-ray14; Mondrian sets on BDCXR-3257).}
\label{sfig:conformal}
\end{figure*}

\begin{figure*}[!t]
\centering
\safefig[0.8\textwidth]{figures/fig_selective_classification_densenet121.pdf}
\caption{Selective classification: risk--coverage curves with the entropy of the seed ensemble alone and with the suppressor's spread
added.}
\label{sfig:selective}
\end{figure*}

\begin{figure*}[!t]
\centering
\safefig[0.9\textwidth]{figures/fig_xai_ig_densenet121.pdf}
\caption{Integrated-Gradients maps of the hybrid model on NIH ChestX-ray14 test images, per input.}
\label{sfig:ig}
\end{figure*}

\begin{figure*}[!t]
\centering
\safefig[0.95\textwidth]{figures/fig_pointing_game_densenet121.pdf}
\caption{Pointing game per NIH box class and per group. (a) Hit rate of each map. (b) Decoder routing: share of each class query's
attention on each stream. (c) Share of the Integrated-Gradients attribution per input. Brackets: number of targets.}
\label{sfig:pointing}
\end{figure*}

\end{document}
